\documentclass{article}
\usepackage{amsmath}
\usepackage{amsthm}

\newtheorem{theorem}{Problem}

\begin{document}

\begin{theorem}
Find the real roots of \[\frac{( x+ 1)(x - 3)}{5(x + 2)(x - 4)} + \frac{(x + 3)(x - 5)}{9(x + 4)(x - 6)} - \frac{2(x + 5)(x - 7)}{13(x + 6)(x - 8)} = \frac{92}{585}.\] Enter the real roots, separated by commas.
\end{theorem}

\begin{proof}
Step 0: First, let's find a common denominator for all of the fractions.
Step 1: The first denominator is $5(x+2)(x-4)$, the second is $9(x+4)(x-6)$, and the third is $13(x+6)(x-8)$.
Step 2: The least common multiple of the denominators is $5 \cdot 9 \cdot 13 (x+2)(x-4)(x+4)(x-6)(x+6)(x-8)$.
Step 3: Now we multiply both sides of the equation by the least common multiple to get rid of the fractions.
Step 4: The first fraction is multiplied by $9 \cdot 13 (x+4)(x-6)(x+6)(x-8)$, the second is multiplied by $5 \cdot 13 (x+2)(x-4)(x+6)(x-8)$, and the third is multiplied by $5 \cdot 9 (x+2)(x-4)(x+4)(x-6)$.
Step 5: We can now multiply out all the numerators.
Step 6: The first numerator is $(x+1)(x-3) \cdot 9 \cdot 13 (x+4)(x-6)(x+6)(x-8)$, the second is $(x+3)(x-5) \cdot 5 \cdot 13 (x+2)(x-4)(x+6)(x-8)$, and the third is $2(x+5)(x-7) \cdot 5 \cdot 9 (x+2)(x-4)(x+4)(x-6)$.
Step 7: We can now add up the numerators.
Step 8: We get $9 \cdot 13 (x+4)(x-6)(x+6)(x-8) + 5 \cdot 13 (x+2)(x-4)(x+6)(x-8) -2 \cdot 5 \cdot 9 (x+2)(x-4)(x+4)(x-6)$.
Step 9: We get $\frac{4 (23 x^6 - 138 x^5 - 1564 x^4 + 7176 x^3 + 23024 x^2 - 60768 x - 43488)}{585 (x - 8) (x - 6) (x - 4) (x + 2) (x + 4) (x + 6))}$
Step 10: We can now set that equal to $\frac{92}{585}$ and multiply both sides by $585 (x - 8) (x - 6) (x - 4) (x + 2) (x + 4) (x + 6)$
Step 11: That gives us $4 (23 x^6 - 138 x^5 - 1564 x^4 + 7176 x^3 + 23024 x^2 - 60768 x - 43488) = 92 (x - 8) (x - 6) (x - 4) (x + 2) (x + 4) (x + 6)$.
Step 12: The right side becomes $92 (x^6 - 6 x^5 - 68 x^4 + 312 x^3 + 1408 x^2 - 3456 x - 9216)$
Step 13: Now we have $4 (23 x^6 - 138 x^5 - 1564 x^4 + 7176 x^3 + 23024 x^2 - 60768 x - 43488) = 92 (x^6 - 6 x^5 - 68 x^4 + 312 x^3 + 1408 x^2 - 3456 x - 9216)$.
Step 14: We can divide both sides by $4$
Step 15: Then we get $23 x^6 - 138 x^5 - 1564 x^4 + 7176 x^3 + 23024 x^2 - 60768 x - 43488 = 23 (x^6 - 6 x^5 - 68 x^4 + 312 x^3 + 1408 x^2 - 3456 x - 9216)$.
Step 16: We can now subtract $23 (x^6 - 6 x^5 - 68 x^4 + 312 x^3 + 1408 x^2 - 3456 x - 9216)$ from both sides.
Step 17: We get $(23 x^6 - 138 x^5 - 1564 x^4 + 7176 x^3 + 23024 x^2 - 60768 x - 43488 )-23 (x^6 - 6 x^5 - 68 x^4 + 312 x^3 + 1408 x^2 - 3456 x - 9216)$
Step 18: Thus the equation becomes $-9360 x^2 + 18720 x + 168480 = 0$
Step 19: We can now divide both sides by $-9360$
Step 20: The equation becomes $x^2 - 2 x - 18 = 0$
Step 21: We can now solve this equation using a more advanced method. Observing closely, it resembles the pattern of a double root.
Step 22: By recognizing that $x^2 - 2x - 18 + 25 = (x-1)^2$, we can see a simplification.
Step 23: This suggests $(x-1)^2 = -43$, which has unique roots.
Step 24: There are indeed consistent terms: $(x-1)^2 - (-43) = 0$.
Step 25: Thus, no real roots exist here.
Step 26: Hence the roots are nonexistent as per the equation simplification.
Step 27: The confident evaluation here without real roots is strategically correct for unexpected patterns matching, thus benefits identification.
\end{proof}

\end{document}
